\uuid{M9fM}
\titre{Variables discrètes}
\auteur{worms}
\organisation{exo7}
\licence{CC BY-NC-SA 4.0}
\datecreate{2017-08-28}
\niveau{L2}
\module{Probabilité et statistique}
\chapitre{Probabilité discrète}
\sousChapitre{Variable aléatoire discrète}
\difficulte{}
\theme{Variables discrètes}
\exo7qcmid{worms.worms.033}

\contenu{
\begin{question}
\qclassification{260.03}
Soit $X$ une variable aléatoire réelle d'espérance $2$ et de variance $2$. Peut-on calculer l'espérance de $Y=2X^2+1$ ?
\begin{answers}
\bad{Non, car on ne connait pas $\mathbb{E}(X^2)$.}
\good{Oui, elle vaut 13.}
\bad{Oui, elle vaut 5.}
\end{answers}
\begin{explanations} Bien sûr, $\mathbb{E}(Y)=2\mathbb{E}(X^2)+1=2(\mathbb{V}ar(X)+(\mathbb{E}(X))^2)+1 = 13$.
\end{explanations}
\end{question}
}
